% =============================================================================
% NOTA AO LEITOR — Ponte ética entre a moldura arquivística e a leitura
% =============================================================================
\cleardoublepage
\thispagestyle{empty}
\chapter*{Nota ao Leitor}
\addcontentsline{toc}{chapter}{Nota ao Leitor}
\markboth{Nota ao Leitor}{}

\begingroup

O volume que você tem em mãos não é apenas uma obra de ficção. Ele se assenta sobre uma ferida histórica real: o Hospital Colônia de Barbacena (1903--1994), onde mais de 60.000 pessoas foram internadas à revelia e cerca de 1.853 corpos foram exumados de valas comuns nas décadas de 1980 e 1990.

\vspace{0.6em}

Joaquim Antônio Correia é um personagem. Mas a dor que ele carrega — a fome, o deslocamento forçado, o confinamento, a violência institucional — não é ficcional. Milhares de brasileiros anônimos viveram essa mesma travessia. O horror que você está prestes a experimentar tem raízes no mundo real.

\vspace{0.6em}

A manipulação que este livro exerce sobre o leitor é uma escolha formal consciente. Ela não existe para enganar, mas para fazer o corpo sentir o que a mente sabe: que a vala coletiva do Colônia não é uma metáfora. O texto a reconstitui como \textit{experiência} — e, ao fazê-lo, cobra de quem lê o preço ético de testemunhar.

\vspace{0.6em}

Este livro é uma obra de horror, mas o horror que ele convoca não deve ser confundido com entretenimento. As violências que atravessam estas páginas — a fome, o confinamento, a tortura psiquiátrica — foram reais para dezenas de milhares de brasileiros anônimos. Se, em algum momento, a excitação estética da leitura se sobrepuser à memória dos que sofreram, {\itshape pause. Reflita.} O sofrimento documentado não existe para ser consumido. Existe para ser testemunhado.

\vspace{0.6em}

{\itshape Você será manipulado. Você será contaminado. Mas, ao final, saberá por quê.}

\vspace{1.2em}

\begin{flushright}
--- O Arquivista\par
Julho de 2025
\end{flushright}

\endgroup
\clearpage
